\section{Variational Argument for Symmetry Selection}
  \label{app:variational}

  \subsection{Setup and Admissible Representations}
    \label{subsec:setup-and-admissible-representations}

    We formalize the symmetry selection principle of Section~6 through an explicit variational
    argument.
    Consider a quasi-two-dimensional system with point group $G = D_{4h}$ and a dominant
    frustration wavevector $\mathbf{Q} = (\pi,\pi)$.
    The gap function $\Delta(\mathbf{k})$ must belong to an irreducible representation (irrep)
    of $G$.
    The four leading irreps are
    \begin{align}
      A_{1g}:&\quad \Delta_s(\mathbf{k}) = \Delta_0, \nonumber\
      B_{1g}:&\quad \Delta_d(\mathbf{k}) = \Delta_0(\cos k_x - \cos k_y), \nonumber\
      B_{2g}:&\quad \Delta_{d'}(\mathbf{k}) = \Delta_0 \sin k_x \sin k_y, \nonumber\
      A_{1g}^{\pm}:&\quad \Delta_{s^\pm}(\mathbf{k}) = \Delta_0(\cos k_x + \cos k_y).
      \label{eq:irreps}
    \end{align}

  \subsection{The Raccordement Cost Functional}
    \label{subsec:the-raccordement-cost-functional}

    The fiber raccordement cost functional of Eq.~(30) takes the explicit form
    \begin{equation}
      C_\textrm{tot}[\Delta] \;=\; \int_\textrm{FS} \frac{d\Omega_{\mathbf{k}}}{(2\pi)^2}
      \Bigl[r_F^2\,\bigl|\Delta(\mathbf{k} + \mathbf{Q}) + \Delta(\mathbf{k})\bigr|^2
      + (1 - r_F^2)\,\bigl|\nabla_{\mathbf{k}}\Delta(\mathbf{k})\bigr|^2\Bigr],
      \label{eq:Ctot}
    \end{equation}
    where $r_F \in [0,1]$ is the frustration ratio of Eq.~(31) and the integral runs over the
    Fermi surface with solid-angle measure $d\Omega_{\mathbf{k}}$.
    The first term penalises configurations whose sign does not alternate under translation by
    $\mathbf{Q}$; the second penalises rapid phase variations along the Fermi surface.

  \subsection{Minimization: $d_{x^2-y^2}$ as the Global Minimum}
    \label{subsec:minimization-as-the-global-minimum}

    \textbf{Strongly staggered regime ($r_F \to 1$).}
    The dominant contribution to $C_\textrm{tot}$ is the staggered mismatch term.
    For the $B_{1g}$ gap one verifies directly that
    \begin{equation}
      \Delta_d(\mathbf{k}+(\pi,\pi))
      = \Delta_0\bigl(\cos(k_x+\pi) - \cos(k_y+\pi)\bigr)
      = -\Delta_0(\cos k_x - \cos k_y) = -\Delta_d(\mathbf{k}),
      \label{eq:dwave_stagger}
    \end{equation}
    so that the staggered mismatch $|\Delta(\mathbf{k}+\mathbf{Q}) + \Delta(\mathbf{k})|^2$
    vanishes identically for $B_{1g}$.
    By contrast, the $s$-wave and $s^\pm$ irreps both contribute
    $4r_F^2|\Delta_0|^2 > 0$.
    Therefore,
    \begin{equation}
      \underset{\Delta \in \mathrm{irreps}(D_{4h})}{\arg\min}\; C_\textrm{tot}[\Delta]
      \;=\; B_{1g} \quad \text{for } r_F \gtrsim r_F^*.
      \label{eq:dwave_min}
    \end{equation}

    \textbf{Weakly staggered regime ($r_F \to 0$).}
    The gradient term dominates.
    On the cuprate or nickelate Fermi surface, the extended $s^\pm$ form factor accumulates a
    smaller gradient cost than the pure $s$-wave because its nodes lie away from the principal
    Fermi arcs.
    For the nickelate multi-band geometry, inter-orbital hybridization further reduces
    $C_\textrm{tot}$ in the $A_{1g}^\pm$ channel relative to $B_{1g}$, yielding
    \begin{equation}
      \underset{\Delta}{\arg\min}\; C_\textrm{tot}[\Delta] \;=\; A_{1g}^\pm \quad
      \text{for } r_F \lesssim r_F^*.
      \label{eq:spm_min}
    \end{equation}

  \subsection{Required Hypotheses}
    \label{subsec:required-hypotheses}

    The variational argument rests on three independently testable assumptions.
    First, the dominant frustration wavevector is $\mathbf{Q} = (\pi,\pi)$; incommensurate
    channels would shift the nodal structure.
    Second, the Fermi surface possesses $D_{4h}$ symmetry; orthorhombic distortions mix
    $B_{1g}$ and $B_{2g}$ and can deflect the nodes from the zone diagonal.
    Third, the cost functional Eq.~\eqref{eq:Ctot} captures the leading symmetry-breaking term;
    higher-order corrections may renormalise $r_F^*$ but do not change the qualitative
    $B_{1g}$ versus $A_{1g}^\pm$ selection.
    All three assumptions are independently testable through neutron scattering, quantum
    oscillations, and gap-symmetry measurements.
